% --- CORRECCION: Forzar interpretacion raw de la entrada ---
\UseRawInputEncoding
\documentclass[aspectratio=169, 10pt, xcolor=table]{beamer}

% --- Configuracion del Tema ---
\usetheme{Madrid}
\usecolortheme{dolphin}

% --- Paquetes necesarios ---
\usepackage[utf8]{inputenc}
\usepackage[spanish]{babel}
\usepackage{amsmath, amssymb}
\usepackage{booktabs}
\usepackage{graphicx}
\usepackage{listings}
\usepackage{tikz}
\usetikzlibrary{arrows.meta, positioning, shapes.geometric}
\usepackage{etoolbox}

% --- Ruta de Imagenes ---
\graphicspath{{./imagenes/}}

% --- Estilo para codigo Python ---
\definecolor{codegreen}{rgb}{0,0.6,0}
\definecolor{codegray}{rgb}{0.5,0.5,0.5}
\definecolor{codepurple}{rgb}{0.58,0,0.82}
\definecolor{backcolour}{rgb}{0.95,0.95,0.92}
\lstset{
    language=Python,
    backgroundcolor=\color{backcolour},
    commentstyle=\color{codegreen},
    keywordstyle=\color{blue},
    numberstyle=\tiny\color{codegray},
    stringstyle=\color{codepurple},
    basicstyle=\ttfamily\scriptsize,
    breaklines=true,
    frame=single,
    showstringspaces=false,
    literate={á}{{\'a}}1 {é}{{\'e}}1 {í}{{\'i}}1 {ó}{{\'o}}1 {ú}{{\'u}}1 {ñ}{{\~n}}1 {ü}{{\"u}}1
}

% --- Pie de pagina personalizado ---
\makeatletter
\setbeamertemplate{footline}{
  \leavevmode%
  \hbox{%
  \begin{beamercolorbox}[wd=.333333\paperwidth,ht=2.25ex,dp=1ex,center]{author in head/foot}%
    \usebeamerfont{author in head/foot}\insertshortauthor
  \end{beamercolorbox}%
  \begin{beamercolorbox}[wd=.333333\paperwidth,ht=2.25ex,dp=1ex,center]{title in head/foot}%
    \usebeamerfont{title in head/foot}Sesion 7
  \end{beamercolorbox}%
  \begin{beamercolorbox}[wd=.333333\paperwidth,ht=2.25ex,dp=1ex,right]{date in head/foot}%
    \usebeamerfont{date in head/foot}CALCULO Y ALGEBRA LINEAL \hfill \insertframenumber/\inserttotalframenumber \hspace*{0.3cm}
  \end{beamercolorbox}}%
  \vskip0pt%
}
\makeatother

% --- Datos de la portada ---
\title[SESION 7 CALCULO]{\textbf{Sesion 7}}
\subtitle{Regla de la Cadena y Derivadas de Orden Superior}
\author[Prof. C. Sanchez]{Carlos Cesar Sanchez Coronel}
\institute{\textbf{CALCULO Y ALGEBRA LINEAL PARA IA}}
\date{}

% --- Eliminar barra de navegacion ---
\setbeamertemplate{navigation symbols}{}

% --- Indice al inicio de cada seccion ---
\AtBeginSection[]{
  \begin{frame}
    \frametitle{Contenido}
    \tableofcontents[currentsection]
  \end{frame}
}

\begin{document}

% --- Portada ---
\begin{frame}
    \titlepage
\end{frame}

% --- Indice general ---
\begin{frame}{Contenido}
    \tableofcontents
\end{frame}

% ============================================================================
% SECCION 1: COMPOSICION DE FUNCIONES
% ============================================================================
\section{Composicion de Funciones}

\begin{frame}{Funcion compuesta}
    \begin{block}{Definicion}
        Dadas dos funciones $f$ y $g$, la composicion $f\circ g$ se define como
        \[
        (f\circ g)(x) = f(g(x))
        \]
        Primero se aplica $g$, luego $f$ al resultado.
    \end{block}
    \begin{exampleblock}{Ejemplo}
        Sean $g(x)=x^2+1$ y $f(u)=\sin u$. Entonces
        \[
        h(x)=f(g(x))=\sin(x^2+1)
        \]
    \end{exampleblock}
    \begin{itemize}
        \item Las redes neuronales son composiciones enormes de funciones.
    \end{itemize}
\end{frame}

% ============================================================================
% SECCION 2: REGLA DE LA CADENA
% ============================================================================
\section{Regla de la Cadena}

\begin{frame}{Formulaciones de la regla de la cadena}
    \begin{columns}[t]
        \column{0.5\textwidth}
        \textbf{Forma de Leibniz}
        \[
        \frac{dy}{dx} = \frac{dy}{du}\cdot\frac{du}{dx}
        \]
        \column{0.5\textwidth}
        \textbf{Forma de Lagrange}
        \[
        (f\circ g)'(x) = f'(g(x))\cdot g'(x)
        \]
    \end{columns}
    \vspace{0.3cm}
    \begin{exampleblock}{Ejemplo}
        Derivar $h(x)=\sin(x^2+1)$.
        \begin{align*}
            u &= x^2+1,\quad h=\sin u \\
            h'(x) &= \cos u \cdot (2x) = 2x\cos(x^2+1)
        \end{align*}
    \end{exampleblock}
\end{frame}

\begin{frame}{Regla de la cadena generalizada}
    Para composiciones multiples, se aplica sucesivamente:
    \[
    \frac{d}{dx}f(g(h(x))) = f'(g(h(x))) \cdot g'(h(x)) \cdot h'(x)
    \]
    \begin{exampleblock}{Ejemplo con tres niveles}
        Derivar $k(x)=e^{\cos(x^2)}$.
        \begin{align*}
            a &= x^2,\quad b=\cos a,\quad k=e^b \\
            k'(x) &= e^b \cdot (-\sin a) \cdot 2x \\
            &= -2x e^{\cos(x^2)} \sin(x^2)
        \end{align*}
    \end{exampleblock}
\end{frame}

% ============================================================================
% SECCION 3: APLICACION EN IA: BACKPROPAGATION
% ============================================================================
\section{Aplicacion en IA: Backpropagation}

\begin{frame}{Red neuronal simple}
    \begin{itemize}
        \item Una entrada $x$, una capa oculta con una neurona (activacion sigmoide) y una salida lineal.
    \end{itemize}
    \begin{align*}
        z &= w_1 x + b_1 \\
        a &= \sigma(z) = \frac{1}{1+e^{-z}} \\
        \hat{y} &= w_2 a + b_2 \\
        L &= \frac{1}{2}(\hat{y}-y)^2 \quad\text{(MSE)}
    \end{align*}
    \begin{center}
        \begin{tikzpicture}[
            neuron/.style={circle, draw, minimum size=0.7cm},
            font=\small
        ]
        \node[neuron] (x) at (0,0) {$x$};
        \node[neuron] (z) at (2,0.5) {$z$};
        \node[neuron] (a) at (4,0) {$a$};
        \node[neuron] (y) at (6,0) {$\hat{y}$};
        \draw[->] (x) -- node[above]{$w_1$} (z);
        \draw[->] (z) -- node[above]{$\sigma$} (a);
        \draw[->] (a) -- node[above]{$w_2$} (y);
        \end{tikzpicture}
    \end{center}
\end{frame}

\begin{frame}{Calculo de gradientes por regla de la cadena}
    \begin{block}{Derivadas parciales}
        \begin{align*}
            \frac{\partial L}{\partial \hat{y}} &= \hat{y}-y \\
            \frac{\partial L}{\partial w_2} &= \frac{\partial L}{\partial \hat{y}}\cdot\frac{\partial \hat{y}}{\partial w_2} = (\hat{y}-y)\,a \\
            \frac{\partial L}{\partial b_2} &= \hat{y}-y
        \end{align*}
        Para $w_1$:
        \[
        \frac{\partial L}{\partial w_1} = \frac{\partial L}{\partial \hat{y}}\cdot\frac{\partial \hat{y}}{\partial a}\cdot\frac{\partial a}{\partial z}\cdot\frac{\partial z}{\partial w_1}
        \]
        donde $\frac{\partial \hat{y}}{\partial a}=w_2$, $\frac{\partial a}{\partial z}=a(1-a)$, $\frac{\partial z}{\partial w_1}=x$.
    \end{block}
\end{frame}

\begin{frame}{Ejemplo numerico paso a paso}
    \begin{exampleblock}{Datos: $x=2$, $w_1=0.5$, $b_1=0.1$, $w_2=1.0$, $b_2=0.0$, $y=5.0$}
        \begin{align*}
            z &= 0.5\cdot2 + 0.1 = 1.1 \\
            a &= \sigma(1.1) = \frac{1}{1+e^{-1.1}}\approx 0.75026 \\
            \hat{y} &= 1.0\cdot0.75026 + 0.0 = 0.75026 \\
            L &= \frac{1}{2}(0.75026-5)^2 \approx 9.031
        \end{align*}
        \begin{align*}
            \frac{\partial L}{\partial \hat{y}} &= -4.24974 \\
            \frac{\partial L}{\partial w_2} &= -4.24974 \times 0.75026 \approx -3.188 \\
            \frac{\partial L}{\partial w_1} &= (-4.24974) \times 1.0 \times 0.75026(1-0.75026) \times 2 \approx -1.593
        \end{align*}
    \end{exampleblock}
\end{frame}

% ============================================================================
% SECCION 4: DERIVADAS DE ORDEN SUPERIOR
% ============================================================================
\section{Derivadas de Orden Superior}

\begin{frame}{Segunda derivada y curvatura}
    \begin{block}{Definicion}
        $f''(x) = \frac{d}{dx}\left(\frac{df}{dx}\right) = \frac{d^2f}{dx^2}$
    \end{block}
    \begin{itemize}
        \item $f''(x) > 0$: funcion convexa (curva hacia arriba).
        \item $f''(x) < 0$: funcion concava (curva hacia abajo).
        \item $f''(x) = 0$: posible punto de inflexion.
    \end{itemize}
    \begin{exampleblock}{Ejemplo}
        $f(x)=x^3-3x$:
        \[
        f'(x)=3x^2-3,\quad f''(x)=6x
        \]
        En $x=1$: $f''(1)=6>0$ (minimo). En $x=-1$: $f''(-1)=-6<0$ (maximo). En $x=0$: $f''(0)=0$ (punto de inflexion).
    \end{exampleblock}
\end{frame}

\begin{frame}{Matriz Hessiana}
    \begin{block}{Para funciones de varias variables}
        La matriz Hessiana $\mathbf{H}$ contiene las segundas derivadas parciales:
        \[
        \mathbf{H}_{ij} = \frac{\partial^2 f}{\partial x_i \partial x_j}
        \]
    \end{block}
    \begin{itemize}
        \item En optimizacion, el Hessiano indica la curvatura de la funcion de perdida.
        \item Metodos de segundo orden (Newton) usan el Hessiano para acelerar convergencia.
        \item Los puntos de silla (comunes en altas dimensiones) tienen Hessiano indefinido.
    \end{itemize}
\end{frame}

% ============================================================================
% SECCION 5: EJEMPLOS CON PYTHON
% ============================================================================
\section{Ejemplos con Python}

\begin{frame}[fragile]{Derivadas de orden superior con SymPy}
    \begin{lstlisting}
import sympy as sp

x = sp.Symbol('x')
f = x**3 * sp.exp(x) - sp.ln(x)

# Primera derivada
f1 = sp.diff(f, x)
print("f'(x) =", f1)

# Segunda derivada
f2 = sp.diff(f, x, 2)
print("f''(x) =", f2)

# Evaluar en un punto
valor = f2.subs(x, 1).evalf()
print("f''(1) =", valor)
    \end{lstlisting}
\end{frame}

\begin{frame}[fragile]{Derivadas parciales y Hessiano simbolico}
    \begin{lstlisting}
import sympy as sp

x, y = sp.symbols('x y')
f = x**2 * y + sp.sin(y)

# Derivadas parciales
df_dx = sp.diff(f, x)
df_dy = sp.diff(f, y)
print("df/dx =", df_dx)
print("df/dy =", df_dy)

# Hessiano
hess = sp.hessian(f, (x, y))
print("Hessiano:\n", hess)
    \end{lstlisting}
\end{frame}

\begin{frame}[fragile]{Backpropagation manual en Python}
    \begin{lstlisting}
import numpy as np

x, y_true = 2.0, 5.0
w1, b1 = 0.5, 0.1
w2, b2 = 1.0, 0.0

# Forward
z = w1 * x + b1
a = 1 / (1 + np.exp(-z))
y_pred = w2 * a + b2
loss = 0.5 * (y_pred - y_true)**2

print(f"Prediccion: {y_pred:.4f}, Loss: {loss:.4f}")

# Backward
dL_dy = y_pred - y_true
dy_da = w2
da_dz = a * (1 - a)
dz_dw1 = x
dz_db1 = 1.0
grad_w1 = dL_dy * dy_da * da_dz * dz_dw1
grad_b1 = dL_dy * dy_da * da_dz * dz_db1
grad_w2 = dL_dy * a
grad_b2 = dL_dy

print("Gradientes:", grad_w1, grad_b1, grad_w2, grad_b2)
    \end{lstlisting}
\end{frame}

\begin{frame}[fragile]{Verificacion con JAX (diferenciacion automatica)}
    \begin{lstlisting}
import jax
import jax.numpy as jnp

def forward(params, x):
    w1, b1, w2, b2 = params
    z = w1 * x + b1
    a = 1 / (1 + jnp.exp(-z))
    y_pred = w2 * a + b2
    return y_pred

def loss(params, x, y_true):
    y_pred = forward(params, x)
    return 0.5 * (y_pred - y_true)**2

params = (0.5, 0.1, 1.0, 0.0)
x_val, y_val = 2.0, 5.0

grad_loss = jax.grad(loss)
grads = grad_loss(params, x_val, y_val)
print("Gradientes con JAX:", grads)
    \end{lstlisting}
\end{frame}

% ============================================================================
% SECCION 6: NOTACION EN PAPERS
% ============================================================================
\section{Notacion en Papers Cientificos}

\begin{frame}{Notacion comun en Deep Learning}
    \begin{itemize}
        \item $\frac{\partial \mathcal{L}}{\partial \mathbf{W}^{(l)}} = \delta^{(l)} (\mathbf{a}^{(l-1)})^{\top}$: gradiente de la perdida respecto a pesos de la capa $l$.
        \item $\delta^{(l)}$: error propagado hacia atras.
        \item $\mathbf{H} = \nabla^2_{\mathbf{w}} \mathcal{L}(\mathbf{w})$: matriz Hessiana.
        \item Actualizacion de Newton: $\mathbf{w}_{t+1} = \mathbf{w}_t - \eta \mathbf{H}^{-1} \nabla \mathcal{L}(\mathbf{w}_t)$.
    \end{itemize}
\end{frame}

% ============================================================================
% SECCION 7: EJERCICIOS PROPUESTOS
% ============================================================================
\section{Ejercicios Propuestos}

\begin{frame}{Ejercicios}
    \begin{enumerate}
        \item Deriva $f(x) = \ln(\sin(e^x))$ usando regla de la cadena paso a paso.
        \item Calcula la segunda derivada de $g(x) = e^{x^2+1}$.
        \item Para la red neuronal del ejemplo, obtén las derivadas parciales si la activacion fuera ReLU en lugar de sigmoide. ¿Que cambios observas?
        \item Implementa en Python una red con dos neuronas en la capa oculta y calcula los gradientes manualmente (usa un bucle sobre las neuronas).
        \item Investiga: ¿que es el "vanishing gradient" y como se relaciona con la regla de la cadena?
    \end{enumerate}
\end{frame}

% ============================================================================
% SECCION 8: RESUMEN
% ============================================================================
\section{Resumen}

\begin{frame}{Lo que hemos aprendido}
    \begin{exampleblock}{Conceptos clave}
        \begin{itemize}
            \item \textbf{Regla de la cadena}: derivada de funciones compuestas, base del backpropagation.
            \item \textbf{Aplicacion en IA}: calculo de gradientes en redes neuronales simples.
            \item \textbf{Derivadas de orden superior}: curvatura, segunda derivada, matriz Hessiana.
            \item \textbf{Optimizacion}: criterios de minimos/maximos y metodos de segundo orden.
            \item \textbf{Implementacion}: SymPy para derivadas simbolicas, NumPy para backprop manual, JAX para diferenciacion automatica.
        \end{itemize}
    \end{exampleblock}
\end{frame}

\begin{frame}{Conexion con futuros cursos}
    \begin{itemize}
        \item \textbf{Deep Learning}: backpropagation en redes profundas.
        \item \textbf{Optimizacion avanzada}: metodos de Newton, quasi-Newton.
        \item \textbf{Teoria del aprendizaje}: analisis de convergencia y puntos criticos.
    \end{itemize}
\end{frame}

% --- Diapositiva final ---
\begin{frame}
    \centering
    \Huge \textbf{Gracias!}

\end{frame}

\end{document}